\section*{الفصل \ref{ch:g3:bones-and-muscles} --- العظام والعضلات}

\begin{solution}{exo:g3:bones-and-muscles:1}
\omterm{def:g3:bones-and-muscles:skeleton}{الجمجمة} تحمي الدماغ؛ \omterm{def:g3:bones-and-muscles:skeleton}{والعمود الفقري} يحفظ \omterm{def:g1:my-body:parts}{البدن} منتصبًا (وهو عمود
قابل للانثناء)؛ \omterm{def:g3:bones-and-muscles:skeleton}{والأضلاع} تكوّن قفصًا واقيًا حول القلب والرئتين؛
\omterm{def:g3:bones-and-muscles:skeleton}{وعظام} \omterm{def:g1:my-body:parts}{الأطراف} تحمل \omterm{def:g1:my-body:parts}{الذراعين} \omterm{def:g1:my-body:parts}{والرجلين}.
\end{solution}

\begin{solution}{exo:g3:bones-and-muscles:2}
نحو $200$. وأطولها \omterm{def:g3:bones-and-muscles:skeleton}{عظم} الفخذ.
\end{solution}

\begin{solution}{exo:g3:bones-and-muscles:3}
لا --- \omterm{def:g3:bones-and-muscles:skeleton}{فالعظم} المنثني \omterm{def:g3:bones-and-muscles:skeleton}{عظم} مكسور. وكل انثناء يحدث عند \omterm{prop:g3:bones-and-muscles:joints}{المفاصل}، حيث
تلتقي عظمتان.
\end{solution}

\begin{solution}{exo:g3:bones-and-muscles:4}
إنه عمود من أكثر من $20$ \omterm{def:g3:bones-and-muscles:skeleton}{عظمة} صغيرة، تستطيع كل واحدة أن تدور
درجات قليلة على جارتيها؛ والانثناءات الصغيرة الكثيرة تتجمع في
التفافة كبيرة لدنة واحدة.
\end{solution}

\begin{solution}{exo:g3:bones-and-muscles:5}
تقصر وتغلظ، فتسحب \omterm{def:g3:bones-and-muscles:skeleton}{العظمة} المتصلة بها. ولا \omterm{def:g3:bones-and-muscles:muscle}{عضلة} تستطيع أن تدفع.
\end{solution}

\begin{solution}{exo:g3:bones-and-muscles:6}
لأن \omterm{def:g3:bones-and-muscles:muscle}{العضلة} تسحب فحسب، فلا تستطيع \omterm{def:g3:bones-and-muscles:muscle}{عضلة} واحدة أن تحرك \omterm{def:g1:my-body:joint}{المفصل} إلا
في اتجاه واحد. فعند الثني: \omterm{def:g3:bones-and-muscles:muscle}{تنقبض} \omterm{def:g3:bones-and-muscles:muscle}{العضلة} الأمامية فتسحب الساعد إلى
أعلى بينما ترتخي الخلفية. وعند المدّ: \omterm{def:g3:bones-and-muscles:muscle}{تنقبض} الخلفية بينما ترتخي
الأمامية.
\end{solution}

\begin{solution}{exo:g3:bones-and-muscles:7}
\omterm{def:g3:bones-and-muscles:skeleton}{العظم} الحي يعيد بناء الكسر بنفسه ببطء؛ والجبيرة لا تفعل إلا أن
تثبّت القطع بينما يعمل. وهذا يدل على أن \omterm{def:g3:bones-and-muscles:skeleton}{العظم} مادة حية يغذّيها
الدم --- لا ميتة كالطباشير.
\end{solution}

\begin{solution}{exo:g3:bones-and-muscles:8}
عند الثني: ينتفخ مقدَّم \omterm{def:g1:my-body:parts}{الذراع} ويقسو (فتلك \omterm{def:g3:bones-and-muscles:muscle}{العضلة} \omterm{def:g3:bones-and-muscles:muscle}{تنقبض}). وعند
المدّ: يلين المقدَّم بينما تشتدّ مؤخرة \omterm{def:g1:my-body:parts}{الذراع} --- والزوج يتبادل
الدورين.
\end{solution}

\begin{solution}{exo:g3:bones-and-muscles:9}
الركبتان مفصلان مبنيان للطيّات الكبيرة، وتحركهما أقوى \omterm{def:g3:bones-and-muscles:muscle}{عضلات} الجسم؛
أما \omterm{def:g3:bones-and-muscles:skeleton}{العمود الفقري} فمبني لانثناءات رقيقة موزعة على \omterm{def:g3:bones-and-muscles:skeleton}{عظام} صغيرة
كثيرة. والرفع بركبتين منثنيتين وظهر مستقيم يعطي العمل الثقيل
\omterm{def:g1:my-body:parts}{للرجلين} ويوفّر \omterm{def:g3:bones-and-muscles:skeleton}{العمود الفقري}.
\end{solution}

\begin{solution}{exo:g3:bones-and-muscles:10}
كل خيط، كما كل \omterm{def:g3:bones-and-muscles:muscle}{عضلة}، لا يستطيع إلا أن يسحب. وجواب صانع الدمى هو
جواب عالم الأحياء: خيطان لكل حركة --- واحد يسحب \omterm{def:g1:my-body:parts}{الذراع} إلى أعلى،
وآخر يسحبها إلى أسفل --- يعملان بالتناوب، تمامًا كأزواج \omterm{def:g3:bones-and-muscles:muscle}{عضلات}
الجسم.
\end{solution}
